\section{Related Work}
\label{sec:related}

\subsection{Sexism and Hate Speech Detection Datasets}
\label{sec:related:datasets}

The EDOS dataset \cite{kirk2023semeval} evaluated throughout this paper was
built for SemEval-2023 Task 10, which split sexism detection into three
subtasks of increasing granularity: a binary sexist/not-sexist label, a
four-category classification of sexist content, and an eleven-vector
fine-grained taxonomy. Most published EDOS work addresses these subtasks by
fine-tuning a pretrained encoder on the training split, the same strategy
this paper's own encoder baseline
(Sec.~\ref{sec:methodology:baselines}) follows: shared-task submissions
ranged from BERT-based classifiers \cite{rifat2023acsmkrhr} to conventional
machine learning pipelines built on hand-engineered features
\cite{panwar2023edos}, the latter close in spirit to the TF-IDF baseline
used here. Where an LLM appears at all in this literature, it is typically
one comparison point among several fine-tuned models, not the object of a
systematic prompt-design study.

Sexism-specific datasets remain comparatively rare outside English. SWSR
\cite{jiang2021swsr} is, to our knowledge, the only comparably scaled
non-English sexism dataset built with its own dedicated lexicon, drawn from
Sina Weibo rather than the Reddit and Gab sources behind EDOS. Its
existence is a reminder that the platform and language choices behind EDOS
are not incidental: what counts as sexist language, and how directly it
gets stated, varies with both. This study inherits EDOS's English,
Reddit/Gab-sourced scope rather than testing across languages or
platforms, a scope limitation the paper's closing sections return to.

\subsection{Prompt Engineering and In-Context Learning}
\label{sec:related:prompting}

Brown et al. \cite{brown2020gpt3} showed that a sufficiently large frozen
language model can pick up a new task from a short natural-language
description and a handful of demonstrations, with no gradient update at
all. That result opened a design space, mapped by Liu et al.'s survey
\cite{liu2021pretrain}, of choices a practitioner makes when turning a
task into a prompt: how to phrase the instruction, whether to supply
demonstrations and how many, whether to ask the model to reason before
answering. The five toggleable components studied here, role, aspects,
context, chain-of-thought, and few-shot
(Sec.~\ref{sec:methodology:prompts}), occupy one corner of that space,
chosen for a single task rather than surveyed in the abstract.

Chain-of-thought prompting \cite{wei2022chain} is the best-known instance
of the reasoning component in that space, proposed for arithmetic and
multi-step logical problems rather than classification. Sexism detection
is not obviously the kind of task chain-of-thought was built for, and the
results here bear that out unevenly: only two of six models
(Sec.~\ref{sec:results:grid}) show a large, consistent benefit from it,
which argues against assuming a chain-of-thought instruction helps by
default on a new task rather than measuring it directly.

Order matters too. Lu et al. \cite{lu2022fantastically} found that
permuting the order of few-shot demonstrations alone can move a model from
near-state-of-the-art to near-random accuracy on standard classification
benchmarks, with no change to which demonstrations are shown. The order
sweep in this paper (Sec.~\ref{sec:methodology:sweeps}) asks a related but
distinct question: not the order of demonstrations within a fixed prompt,
but the order of entire prompt components (context, chain-of-thought,
aspects, few-shot) relative to each other. The effect is far more
contained than Lu et al.'s general finding: reordering shifts F1 for only
one of six models (Sec.~\ref{sec:results:sensitivity}), suggesting
component-order sensitivity is not a universal property of these models
the way demonstration-order sensitivity appears to be.

\subsection{Large Language Models for Hate Speech and Content Moderation}
\label{sec:related:llm-moderation}

The finding closest in spirit to this paper's own baseline comparison
(Sec.~\ref{sec:results:baselines}) comes from Plaza-del-Arco et al.
\cite{plazadelarco2023label}, who ran several LLMs zero-shot across a
range of text classification tasks and reported that none of them rivaled
even simple supervised baselines. This paper reaches a similar conclusion
by a different route: an 18-variant grid of systematically combined
prompt components per model, rather than zero-shot prompting alone, so
the gap to a linear classifier persists after substantial prompt
engineering, not only in its absence. A comparable pattern shows up in
hate speech detection specifically. Bernal-Beltr\'an et al.
\cite{bernalbeltran2025spanish} compared fine-tuned Spanish encoders
against zero- and few-shot LLM prompting on a hate speech corpus and found
the fine-tuned encoders consistently ahead, with few-shot prompting
narrowing but not closing the gap. Neither study applies the statistical
significance testing or the multi-model breadth used here
(Sec.~\ref{sec:methodology:stats}), but both point the same direction:
for classification tasks like these, a general-purpose LLM's prompting
still lags a model trained on the task's own labels.

Safety alignment introduces a failure mode distinct from accuracy.
R\"ottger et al.'s XSTest \cite{rottger2023xstest} is a purpose-built,
adversarial suite of safe-sounding prompts designed to trigger unwarranted
refusals, and it shows the phenomenon exists in several production models.
The llama3.1 refusal behavior documented in this paper
(Sec.~\ref{sec:results:llama-refusal}) was not sought out through an
adversarial test suite; it surfaced as a byproduct of ordinary sexism
classification once role and aspect framing were dropped from the prompt,
on genuinely sexist and non-sexist text rather than deliberately
borderline phrasing. It is a naturally occurring instance of the same
over-refusal phenomenon XSTest was built to probe, arising from an
ordinary classification pipeline rather than a red-teaming exercise.

\subsection{Fine-Tuned Baselines: Encoders and Parameter-Efficient
Adaptation}
\label{sec:related:baselines}

The encoder baseline in this study
(Sec.~\ref{sec:methodology:baselines}) fine-tunes RoBERTa
\cite{liu2019roberta}, whose central finding, that BERT's original
training recipe left real performance on the table, shows that a
fine-tuning baseline's strength depends on training choices as much as
architecture. The QLoRA baseline builds on two more recent ideas: LoRA
\cite{hu2021lora}, which adapts a frozen model by training a small pair of
low-rank matrices injected into its attention layers instead of updating
all its weights, and QLoRA \cite{dettmers2023qlora}, which quantizes the
frozen base weights to 4 bits so the same adaptation fits on a single GPU
even for a multi-billion-parameter model. Together they make it practical
to fine-tune all six of this paper's prompted models rather than only the
smallest, which is what makes the QLoRA-versus-prompting comparison in
Sec.~\ref{sec:results:baselines} possible at this scale.
